\documentclass[a4paper,11pt]{article}
\usepackage[utf8]{inputenc}
\usepackage[T1]{fontenc}
\usepackage[english]{babel}
\usepackage[left=2.5cm,right=2.5cm,top=2cm,bottom=2cm]{geometry}
\usepackage{hyperref}
\usepackage{xcolor}
\usepackage{microtype}
\usepackage{lmodern}

\definecolor{primary}{RGB}{210, 0, 28}
\hypersetup{colorlinks=true, urlcolor=primary}

\begin{document}
\noindent \textbf{Lo\"ic Aron MBASSI EWOLO} \\
ENSEA -- Double Dipl\^ome Ing\'enieur \\
\href{mailto:loicaron.mbassiewolo@ensea.fr}{loicaron.mbassiewolo@ensea.fr} \quad (+33) 7 59 52 33 98 \\
\href{https://github.com/Nameless0l}{github.com/Nameless0l} \quad \href{https://mbassiloic.tech}{mbassiloic.tech}

\vspace{1em}
\noindent Cergy, le \today

\vspace{1em}
\noindent \textbf{Objet : Candidature -- Alternance Cybersécurité Syst\`emes Critiques (TotalEnergies)}

\vspace{1.5em}
\noindent Madame, Monsieur,

\vspace{0.8em}
\noindent En tant que stagiaire au laboratoire INRIA Saclay (projet PRIMO), je mène une analyse automatisée de risques de sécurité dans 500+ applications mobiles : identification de vecteurs d'attaque via SDKs tiers non sécurisés, permissions excessives, données sensibles (localisation) exposées. Double diplôme ENSEA/Polytechnique Yaoundé en sécurité des systèmes, cryptographie et mathématiques appliquées. TotalEnergies, acteur majeur de la transition énergétique, expose à des défis critiques de cybersécurité OT/IT sur des infrastructures mondiales nécessitant conformité réglementaire stricte et résilience.

\vspace{0.8em}
\noindent J'ai structuré dans le projet PRIMO un pipeline complet d'analyse de risques : instrumentation binaire d'applications, extraction et analyse de 500+ SDKs embarqués, identification de vecteurs d'attaque spécifiques (localisation, données personnelles, connexions non chiffrées). Un scoring automatisé de risques via NLP (BERT, Transformers) sur politiques de confidentialité complète l'analyse statique. J'ai quantifié ces risques dans un ranking interprétable pour équipes de sécurité et décideurs. Cette approche -- combinant reverse engineering, NLP, et gouvernance sécurité -- reproduit les défis de détection et de prévention que TotalEnergies opère sur ses infrastructures critiques OT/IT.

\vspace{0.8em}
\noindent En parallèle, j'ai conçu une architecture de sécurité avancée dans la plateforme Snappy : implémentation du protocole Signal (Perfect Forward Secrecy, double ratchet, AES-256-GCM), chiffrement bout-en-bout, RGPD by design (droit à l'oubli cryptographique, audit logs SIEM), segmentation données critiques. Cette expérience en protocoles cryptographiques et en infrastructure sécurisée démontre une maturité pour adresser les enjeux de sécurité des systèmes critiques. Mes travaux complémentaires en sécurité fintech (CSC, secteur bancaire, conformité PCI-DSS) et en IA détection (Women Safe) enrichissent ce profil. Je suis disponible dès septembre 2026.

\vspace{0.8em}
\noindent TotalEnergies, par ses enjeux de transition énergétique et ses infrastructures mondialement critiques, offre un contexte exceptionnel pour consolider ces compétences en cybersécurité OT/IT et conformité réglementaire. Je serais ravi de discuter comment mon profil peut contribuer à vos initiatives de sécurité et de résilience numériques.

\vspace{1.5em}
\noindent Veuillez agr\'eer, Madame, Monsieur, l'expression de mes salutations distingu\'ees.
\vspace{2em}
\noindent \textbf{Lo\"ic Aron MBASSI EWOLO}

\end{document}
